\fig{em5:polar-sector}{径向、切向速度与半径扫过的微小面积}{
\begin{tikzpicture}[scale=1.05]
\fill[blue!10](0,0)--(15:2.4)arc(15:32:2.4)--cycle;
\draw[world](0,0)--(15:3.3) (0,0)--(32:3.3);
\draw[world](15:2.4)arc(15:32:2.4);
\fill(0,0)circle(1.5pt)node[below left]{$O$};
\coordinate(P)at(15:2.4);
\fill(P)circle(2pt)node[below]{$P$};
\draw[flow,blue](P)--($(P)+(.97,.26)$)node[right]{$v_r=\dot r$};
\draw[flow,red](P)--($(P)+(-.31,1.16)$)node[above]{$v_\theta=r\dot\theta$};
\node[below]at(15:1.3){$r$};
\draw[world](15:.8)arc(15:32:.8);
\node[right]at(24:1){$\dd\theta$};
\node at(24:1.8){$\dd A$};
\end{tikzpicture}}

图\ref{em5:polar-sector}把速度分为沿半径与垂直半径的两部分。径向分量不产生关于$O$的角动量，也不扫出一阶面积；微小扇形面积为$\dd A=\tfrac12r^2\dd\theta$。所以$L=mr v_\theta$与$\dot A=L/(2m)$来自同一几何关系。半径本身变化时也可在当前$r$处作这一瞬时分解，不能把合速度直接当作切向速度。
